\section{自动微分与深度学习框架}

\subsection{前向传播与反向传播的模块化}

反向传播的美妙之处在于其\textbf{模块化}：每个操作只需要实现两个接口：

\begin{enumerate}
    \item \textbf{forward}：给定输入，计算输出（并保存中间值供反向使用）
    \item \textbf{backward}：给定上游梯度，计算对每个输入的下游梯度
\end{enumerate}

\begin{figure}[H]
\centering
\includegraphics[width=0.95\textwidth]{slides-images/slide-069.jpg}
\caption{模块化反向传播：每个操作是一个独立模块，实现 forward 和 backward 接口\protect\footnotemark}
\end{figure}
\footnotetext{Slides 第70页。}

\begin{importantbox}{前向传播保存的值在反向传播中使用}
前向传播时，每个节点不仅需要计算输出，还需要\textbf{缓存}计算中间值。例如：
\begin{itemize}
    \item 乘法 $z = xy$：反向时需要 $x$ 和 $y$ 的值来计算 $\frac{\partial z}{\partial x} = y$ 和 $\frac{\partial z}{\partial y} = x$
    \item ReLU $y = \max(0, x)$：反向时需要知道 $x$ 是否大于 0
    \item Sigmoid $y = \sigma(x)$：反向时可以用 $y$ 来计算 $y(1-y)$
\end{itemize}
这就是为什么训练神经网络比推理需要更多内存——需要保存所有中间激活值。
\end{importantbox}

\subsection{现代深度学习框架}

PyTorch、TensorFlow 等框架的核心能力就是\textbf{自动微分}（Automatic Differentiation）：用户只需定义前向计算，框架自动构建计算图并执行反向传播。

\begin{figure}[H]
\centering
\includegraphics[width=0.95\textwidth]{slides-images/slide-072.jpg}
\caption{自动微分：框架自动追踪前向计算并生成反向传播代码\protect\footnotemark}
\end{figure}
\footnotetext{Slides 第73页。}

\begin{knowledgebox}{静态图 vs 动态图}
\begin{itemize}
    \item \textbf{静态图}（如早期 TensorFlow）：先定义整个计算图，然后执行。优点是可以做全局优化
    \item \textbf{动态图}（如 PyTorch）：逐行执行代码时动态构建计算图。优点是更灵活，更容易调试
\end{itemize}
现代框架（包括 TensorFlow 2.x）已经普遍支持动态图模式，PyTorch 是目前学术研究中最流行的框架。
\end{knowledgebox}

\subsection{本章小结}

反向传播的模块化设计使得自动微分成为可能。现代深度学习框架让用户只需关注前向计算，梯度计算完全自动化。这极大降低了构建和实验新神经网络架构的门槛。

%% ========== 总结与延伸 ==========

